\documentclass[12pt, a4paper]{article}

% --- PREAMBLE: Document Setup ---

% 1. Page Layout: Sets 1-inch margins on all sides.
\usepackage[a4paper, margin=1in]{geometry}

% 2. Font & Language (for pdflatex compatibility)
\usepackage[utf8]{inputenc}     % Handles input text encoding.
\usepackage[T1]{fontenc}        % Handles font encoding for output.
\usepackage{mathptmx}           % Sets the main font to a Times Roman clone.
\usepackage[english]{babel}     % Provides language support.

% Extra packages
\usepackage{graphicx, caption, subcaption}

\usepackage[utf8]{inputenc}
\usepackage{booktabs} % For professional quality lines
\usepackage{siunitx}  % For alignment of numbers and percentages
\usepackage{threeparttable}

\usepackage{multicol} % Enables two-column layout
\usepackage{caption} % Allows for \captionof outside of float environments

% Add to the preamble (after your other \usepackage lines)
\setlength{\fboxsep}{0pt}  % no padding: border touches the image
\setlength{\fboxrule}{3pt} % border thickness (adjust: 0.8pt, 1pt, 1.5pt, ...)
\newcommand{\framedinclude}[2][]{\fbox{\includegraphics[#1]{#2}}}

% 3. Spacing: Sets the entire document to be double-spaced.
\usepackage{setspace}
\doublespacing

% --- DOCUMENT INFORMATION ---
% Define your paper's details here for easy reuse.
% We use \newcommand to define custom commands for the title and author.
\newcommand{\thetitle}{Placeholder}
\newcommand{\theauthor}{Elvis N. Trujillo Chairez}
\newcommand{\affiliation}{University of Houston}
\newcommand{\coursename}{BIOL 4324: Bioinformatices}
\newcommand{\instructor}{Adam Stuckert, PhD}
\newcommand{\duedate}{\today} % Automatically uses today's date, or you can type a specific one.

% We still pass this information to the standard \title and \author commands
% for metadata purposes, which is good practice.
\title{\thetitle}
\author{\theauthor}

% --- DOCUMENT START ---
\begin{document}

% =========================================================================
% --- TITLE PAGE ---
% The 'titlepage' environment creates a separate, unnumbered title page.
% =========================================================================
\begin{titlepage}
    \centering
    \vspace*{3cm} % Adds space from the top of the page.

    % Title of the paper (bold and larger font).
    \textbf{\huge \thetitle}

    \vspace{2cm} % Adds vertical space between the title and author details.

    % Author and affiliation details.
    \large
    \theauthor \\
    \vspace{0.5cm}
    \affiliation \\
    \vspace{0.5cm}
    \coursename \\
    \vspace{0.5cm}
    \instructor \\
    \vspace{0.5cm}
    \duedate

    \vfill % Centers the block of text vertically on the page.
\end{titlepage}


% =========================================================================
% --- ABSTRACT ---
% Starts on a new page.
% =========================================================================
\newpage
\begin{center}
    \textbf{Abstract}
\end{center}
% The first line of the abstract is not indented.
\par

\newpage

% =========================================================================
% --- MAIN BODY ---
% Starts on a new page with the title repeated at the top.
% =========================================================================
\begin{center}
    \textbf{\large Introduction}
\end{center}
\par

\newpage

\begin{center}
    \textbf{\large Material \& Methodology}
\end{center}
\par % This command ensures the text below starts as a new, indented paragraph.

\begin{center}
    \textbf{\large Results}
\end{center}
\par 

\newpage 

\begin{center}
    \textbf{\large Discussion}
\end{center}
\par

% =========================================================================
% --- REFERENCES ---
% Starts on a new page with a centered, bold heading.
% =========================================================================
\newpage
\begin{center}
    \textbf{References}
\end{center}

% Each reference entry is a separate paragraph with a hanging indent.
% \hangindent sets the indentation for all lines of the paragraph except the first.
% \noindent ensures the first line starts at the left margin.

\par\hangindent=0.5in
\noindent

\par\hangindent=0.5in
\noindent

\par\hangindent=0.5in
\noindent

\par\hangindent=0.5in
\noindent

\par\hangindent=0.5in
\noindent

\par\hangindent=0.5in
\noindent

\par\hangindent=0.5in
\noindent

\end{document}